% !TEX root = ../jlt.tex
\section{Featherweight Jlt typing}
\label{app:fjlt:typing}

Figures~\ref{fig:fjlt-entailment} and~\ref{fig:fjlt-typing} say which FJlt programs are well formed, the domain of $\llbracket\cdot\rrbracket$ in Figure~\ref{fig:arrows} and Section~\ref{sec:translation}.
As Section~\ref{sec:fjlt:syntax}'s carrier sets keep \RuleOne{}, their rules hold a program to \RuleTwo{}, through the definition check and resolution, and to \RuleThree{}, through well-founded declarations and ranges (Section~\ref{sec:posture}).
There are three judgments.
$\Delta \Vdash Q$ says that the constraint $Q$ is \emph{entailed} in the context $\Delta$, which holds type variables and assumed constraints.
$\Delta; \Gamma \vdash e : \tau$ types a term, $\Gamma$ holding the term variables with their types.
$\vdash D$ says a declaration is well formed.
A program is well formed when its declarations are and its closed term has a type.
The compiler holds a program to most of this.
Four things rest on the programmer's discipline instead: a generic body held to its constraints (Section~\ref{sec:limits}), a call at a type variable held to resolution, a converted result held to its declared type, and a range held to its conditions.

\begin{figure}[htbp]
\[
\begin{array}{@{}l@{\quad}l@{}}
\textsc{(E-Assume)} & \dfrac{Q \in \Delta}{\Delta \Vdash Q}
\\[12pt]
\textsc{(E-Law)} & \dfrac{\varnothing; \varnothing \vdash e : \texttt{bool} \qquad e \Downarrow \texttt{true}}{\Delta \Vdash e}
\\[12pt]
\textsc{(E-Law-Assumed)} & \dfrac{\texttt{concept } C \texttt{ = } \overline{A} \qquad C\texttt{<}\overline{\tau}\texttt{>} \in \Delta \qquad e = A_i[\overline{\tau}/\overline{\alpha}] \text{ a law}}{\Delta \Vdash e}
\\[12pt]
\textsc{(E-Shape)} & \dfrac{\begin{array}{c} \texttt{concept } C \texttt{ = } \overline{A} \qquad C\texttt{<}\overline{\tau}\texttt{>} \notin \Delta \\ \Delta \Vdash A_i[\overline{\tau}/\overline{\alpha}] \text{ for each } i \end{array}}{\Delta \Vdash C\texttt{<}\overline{\tau}\texttt{>}}
\\[12pt]
\textsc{(E-Op)} & \dfrac{\Delta; \overline{x : \tau_x} \vdash e : \tau' \qquad \tau' \mathrel{S} \tau}{\Delta \Vdash \texttt{requires (}\overline{\texttt{const }\tau_x\; x}\texttt{) \{ \{ }e\texttt{ \} -> }S\texttt{<}\tau\texttt{>; \}}}
\\[12pt]
\textsc{(E-Not)} & \dfrac{\overline{\tau} \text{ concrete} \qquad \texttt{!}C\texttt{<}\overline{\tau}\texttt{>} \notin \Delta \qquad \Delta \not\Vdash C\texttt{<}\overline{\tau}\texttt{>}}{\Delta \Vdash \texttt{!}C\texttt{<}\overline{\tau}\texttt{>}}
\end{array}
\]
\caption{The rules of entailment, $\Delta \Vdash Q$.
Six rules: an assumption, two for laws, a shape, a requirement and a negation.}
\label{fig:fjlt-entailment}
\end{figure}

\paragraph{Entailment and laws.}
\textsc{E-Law} is where a law is a value: entailment evaluates it.
A shape with a law is then satisfied only at instances whose certificate holds.
Laws are closed terms of base type, which keeps typing and evaluation from chasing each other.
A law is evaluated where a record is built, at concrete types.
At a type variable it is entailed only through an assumed shape that lists it, by \textsc{E-Law-Assumed}, as a call is through \textsc{T-Req-Assumed}.
A negated constraint \cppinline{!C<T>} carries no evidence.
\textsc{E-Not} holds at concrete types where $C$ does not, and the well-founded order of the declarations makes its negative premise decidable.
The side condition of \textsc{E-Shape} makes entailment syntax-directed: an assumed shape is its record variable and is never rebuilt field by field.
Theorem~\ref{thm:coherence} needs exactly that.

\begin{figure}[htbp]
\[
\begin{array}{@{}l@{\quad}l@{}}
\multicolumn{2}{@{}l}{\text{\textit{Calls}}}
\\[7pt]
\textsc{(T-Op)} & \dfrac{\begin{array}{c} O\texttt{<}\overline{\tau}\texttt{>} : \overline{\sigma} \to \sigma \text{ (the standard's signature)} \\ \Delta; \Gamma \vdash \overline{e : \sigma} \end{array}}{\Delta; \Gamma \vdash O\texttt{<}\overline{\tau}\texttt{>\{\}(}\overline{e}\texttt{)} : \sigma}
\\[7pt]
\textsc{(T-Conv)} & \dfrac{\Delta; \Gamma \vdash e : \tau' \qquad \tau' \text{ converts to } \tau}{\Delta; \Gamma \vdash \tau\texttt{(}e\texttt{)} : \tau}
\\[7pt]
\textsc{(T-App)} & \dfrac{\begin{array}{c} \Delta; \Gamma \vdash e : O\texttt{<}\overline{\tau}\texttt{>} \qquad O\text{'s call operator is } \forall\overline{\beta}.\,\overline{Q} \Rightarrow \overline{\tau_x} \to \tau_r \\ \Delta \Vdash \overline{Q}[\overline{\tau}/\overline{\alpha}][\overline{\sigma}/\overline{\beta}] \qquad \Delta; \Gamma \vdash \overline{e' : \tau_x[\overline{\tau}/\overline{\alpha}][\overline{\sigma}/\overline{\beta}]} \end{array}}{\Delta; \Gamma \vdash e\texttt{(}\overline{e'}\texttt{)} : \tau_r[\overline{\tau}/\overline{\alpha}][\overline{\sigma}/\overline{\beta}]}
\\[7pt]
\textsc{(T-Call)} & \dfrac{\begin{array}{c} f\texttt{<}\overline{\sigma}\texttt{>} \text{ resolves to } f_i : \forall\overline{\alpha}.\,\overline{Q} \Rightarrow \overline{\tau_x} \to \tau \\ \Delta \Vdash \overline{Q}[\overline{\sigma}/\overline{\alpha}] \qquad \Delta; \Gamma \vdash \overline{e : \tau_x[\overline{\sigma}/\overline{\alpha}]} \end{array}}{\Delta; \Gamma \vdash f\texttt{<}\overline{\sigma}\texttt{>(}\overline{e}\texttt{)} : \tau[\overline{\sigma}/\overline{\alpha}]}
\\[7pt]
\textsc{(T-Trait)} & \dfrac{\begin{array}{c} f\texttt{<}\overline{\sigma}\texttt{>} \text{ resolves to a declaration of type } \overline{\tau_x} \to \tau \\ \Delta; \Gamma \vdash \overline{e : \tau_x} \end{array}}{\Delta; \Gamma \vdash f\texttt{<}\overline{\sigma}\texttt{>(}\overline{e}\texttt{)} : \tau}
\\[7pt]
\multicolumn{2}{@{}l}{\text{\textit{The definition check}}}
\\[7pt]
\textsc{(T-Req-Assumed)} & \dfrac{\begin{array}{c} C\texttt{<}\overline{\tau}\texttt{>} \in \Delta \\ \texttt{\{ }e_r\texttt{ \} -> }S\texttt{<}\tau\texttt{>} \text{ is a requirement of } C \text{ with parameters } \overline{a : \tau_a} \\ \Delta; \Gamma \vdash \overline{e : \tau_a} \end{array}}{\Delta; \Gamma \vdash \tau^{S}\texttt{(}e_r[\overline{e}/\overline{a}]\texttt{)} : \tau}
\\[7pt]
\textsc{(D-Generic)} & \dfrac{\overline{\alpha}, \overline{Q};\; \overline{x : \tau_x} \vdash e : \tau \qquad \rho \in \{\tau, \texttt{auto}\}}{\begin{array}{c}\vdash \texttt{template <typename }\overline{\alpha}\texttt{> requires } \overline{Q} \\ \texttt{constexpr } \rho\; f\texttt{(}\overline{\tau_x\; x}\texttt{) \{ return } e\texttt{; \}}\end{array}}
\\[7pt]
\textsc{(D-Const)} & \dfrac{\varnothing; \varnothing \vdash e : \tau \qquad \rho \in \{\tau, \texttt{auto}\} \qquad e \text{ a constant initialiser}}{\vdash \texttt{constexpr } \rho\; x \texttt{ = } e\texttt{;}}
\\[7pt]
\textsc{(T-Const)} & \dfrac{\vdash \texttt{constexpr } \rho\; x \texttt{ = } e\texttt{;} \text{ declared before, at } \tau \qquad x \notin \Gamma}{\Delta; \Gamma \vdash x : \tau}
\\[7pt]
\multicolumn{2}{@{}l}{\text{\textit{Recursion, from the adopted library}}}
\\[7pt]
\textsc{(T-Fold)} & \dfrac{\begin{array}{c} \Delta; \Gamma \vdash r : R\texttt{<}\tau\texttt{>} \text{, } R \in \mathcal{R} \\ \Delta; \Gamma \vdash a : \sigma \qquad \Delta; \Gamma \vdash f : \sigma \times \tau \to \sigma \end{array}}{\Delta; \Gamma \vdash \texttt{fold\_left(}r\texttt{, }a\texttt{, }f\texttt{)} : \sigma}
\\[7pt]
\textsc{(T-Fold-R)} & \dfrac{\begin{array}{c} \mathit{IsRange}\texttt{<}R\texttt{, }\tau\texttt{>} \in \Delta \qquad \Delta; \Gamma \vdash r : R \\ \Delta; \Gamma \vdash a : \sigma \qquad \Delta; \Gamma \vdash f : \sigma \times \tau \to \sigma \end{array}}{\Delta; \Gamma \vdash \texttt{fold\_left(}r\texttt{, }a\texttt{, }f\texttt{)} : \sigma}
\\[7pt]
\textsc{(T-Cata)} & \dfrac{\begin{array}{c} \Delta; \Gamma \vdash t : \mu F \\ \Delta; \Gamma \vdash \mathit{alg} : \mu F \times \sigma^{k} \to \sigma \text{ at each node of arity } k \end{array}}{\Delta; \Gamma \vdash \texttt{cata(}\mathit{alg}\texttt{, }t\texttt{)} : \sigma}
\\[7pt]
\textsc{(T-Node)} & \dfrac{\Delta; \Gamma \vdash t : O\texttt{<}\overline{\tau}\texttt{>} \qquad O \text{ a node of } F}{\Delta; \Gamma \vdash t : \mu F}
\end{array}
\]
\caption{The rules for terms, $\Delta; \Gamma \vdash e : \tau$, and for generic functions and constants, $\vdash D$.
They come in three groups: calls, the definition check and recursion.
Omitted as textbook: \textsc{T-Var}, \textsc{T-Construct}, \textsc{T-Field}, \textsc{T-Inj}, \textsc{T-Case}, and the rules for constants, pairs, boxes, conditionals and the arithmetic of base types.}
\label{fig:fjlt-typing}
\end{figure}

\paragraph{Resolution.}
A trait call $f\texttt{<}\overline{\tau}\texttt{>}$ at concrete types \emph{resolves} to a declaration of $f$ in two steps, as \Cpp{} resolves an overloaded template.
The first step takes the most specific of the generic functions whose head matches $\overline{\tau}$ and whose context is entailed, the deleted primary among them.
If that is the primary, the second takes its explicit specialisation at $\overline{\tau}$, if there is one.
The deleted primary never resolves, which is how an operation nobody registered fails a shape.
Well-formedness requires resolution to be a function, in two cases.
At concrete types, no two specialisations share a head.
Among the generic functions whose head matches and whose context is entailed, there is a unique most specific one, by head and then by context, as \Cpp{} orders constrained templates; \Cpp{} rejects the program where there is none.
At a type variable, in a body with context $\Delta$, a call resolves only when its choice survives every instantiation of the body.
So every candidate that would be more specific than the chosen one and whose head unifies with the call's types, explicit specialisations included, is excluded by a constraint $\Delta$ decides, entailing it or its negation.
That refines GHC's rule for overlapping instances~\cite{ghc2024overlapping}, which looks at instance heads only and refuses to commit to an instance the types do not yet decide.
\Cpp{} does not hold it.
A call it would refuse is a body that varies with a type through a constraint it did not declare, a breach of Rule~2, which the repository's linter reports where it can see one (Section~\ref{sec:limits}).

\paragraph{Conversions.}
\textsc{T-Conv} types the conversion between carriers that \Cpp{} inserts implicitly and the kernel term writes, $\tau\texttt{(}e\texttt{)}$.
It is total by the condition on conversions of Appendix~\ref{app:fjlt:closure}.
\textsc{E-Op} reads $S$ as equality for \cppinline{std::same_as} and as conversion for \cppinline{std::convertible_to}.
Write $\tau^{S}\texttt{(}e\texttt{)}$ for $e$ when $S$ is \cppinline{std::same_as}, and for $\tau\texttt{(}e\texttt{)}$ when it is \cppinline{std::convertible_to}.
\textsc{T-Req-Assumed} then types $\tau^{S}$ of the pattern, and a use of the requirement keeps its declared type when \textsc{S-Call} substitutes the instance (\textsc{S-Conv}).
A body may therefore use a converted result only at the declared $\tau$.
Using the unconverted result, say to compare two of them at their own type, uses an operation the constraint did not name.
The typing refuses it, \Cpp{} accepts it, and the repository's linter reports it where it can see it.
The library needs the conversion form in three places: where a requirement's expression is an lvalue of the wanted type (\cppinline{const T&} is not \cppinline{same_as<T>}; \cppinline{IsPointed}'s unit), where a step returns a proxy for the carrier (\cppinline{HasNNOStep}'s successor), and where an object is presented by another type (\cppinline{Presents}).

\paragraph{The definition check.}
\textsc{D-Generic} types a generic body under its constraints, and \textsc{T-Req-Assumed} is what makes that a check.
\textsc{D-Op} is the same with a function object's fields among the variables, and \textsc{D-Instance} the same with an empty context.
Inside a constrained body, a call on a constrained type is typed only as an instance of a requirement of a constraint in $\Delta$.
\cppinline{Op{}(x, x)} is the pattern \cppinline{{ Op{}(a, b) }} at $a, b := x$.
\cppinline{unit<T>(Op{})} is the pattern \cppinline{{ unit<T>(Op{}) }}, whose parameters are the variables of the \cppinline{requires} clause free in it, here none.
A trait is reachable the same way.
A variable of a fieldless function-object type denotes that type's one value, so \cppinline{op} of type \cppinline{Op} is \cppinline{Op{}}, which is how Listing~\ref{lst:rule3}'s \cppinline{concat} applies its \cppinline{op}.
The requirements of a shape, in \textsc{T-Req-Assumed} and \textsc{E-Law-Assumed}, are its own and, transitively, those of its nested shapes.
An assumed \cppinline{IsCommutativeMonoid<T, Op>} thus yields the operations and the associativity law of \cppinline{IsMonoid<T, Op>}.
More than one assumed requirement may match.
The rules then take the leftmost shape of the \cppinline{requires} clause and, within it, the first matching requirement in declaration order, so both rules are functions of the body.
Once instantiated, the candidates denote one operation, by Lemma~\ref{lem:resolve}.
A premise $f : \sigma \times \tau \to \sigma$ in \textsc{T-Fold} and \textsc{T-Case} says that $f$ is a function name or a function object.
Its application to arguments of those types is typed at $\sigma$ by \textsc{T-Call}, \textsc{T-Trait}, \textsc{T-Op}, \textsc{T-App} or, at a type variable, \textsc{T-Req-Assumed}.

\paragraph{Operators.}
An operator expression $\ominus e$ or $e_1 \oplus e_2$ is resolved as \Cpp{} resolves it~\cite[{\texttt{over.match.oper}}]{cpp23standard}, the built-in operator a candidate beside the declared functions \texttt{operator}$\oplus$, and the only one when no operand is of class or enumeration type.
Where the choice is a primitive of $\Sigma$, the expression is typed by the primitive's signature, the rule for the arithmetic of base types; otherwise it is the call \texttt{operator}$\oplus$\texttt{(}$e_1$\texttt{, }$e_2$\texttt{)}, typed by \textsc{T-Call} or \textsc{T-Trait}.
Where the operands are of different base types, the primitive acts after the conversions the standard specifies for that operator: the usual arithmetic conversions for arithmetic and comparison\cite[{\texttt{expr.arith.conv}}]{cpp23standard}, and for a shift the promotion of each operand on its own, the result of the left operand's type\cite[{\texttt{expr.shift}}]{cpp23standard}.
Each such conversion is the grammar's $\tau(e)$, placed by the compiler, as Listing~\ref{lst:rule1}'s \cppinline{id_u(42uz) == 42}, a \cppinline{size_t} against an \cppinline{int}.
A declared \texttt{operator}$\oplus$ needs no common type, since its signature is its own: in Listing~\ref{lst:rule3}'s \cppinline{v * a}, a matrix acting on a vector, the operands' types differ and nothing converts.
The candidate selected may be a rewritten or a reversed one\cite[{\texttt{over.match.oper}}]{cpp23standard}: then the expression is the call under the standard's rewrite, \cppinline{a != b} as \cppinline{!(a == b)}, a reversed candidate with its arguments swapped, \cppinline{a < b} as \cppinline{(a <=> b) < 0}, and typing, translation and evaluation read the rewritten call.

\paragraph{Deduced types and constants.}
Where a declaration's result type $\rho$ is \cppinline{auto}, its $\tau$ is the type of the body in the premise, and \textsc{T-Call}, \textsc{T-Trait} and \textsc{T-App}'s $\tau_r$ read that type.
This is \Cpp{}'s return type deduction~\cite[{\texttt{dcl.spec.auto}}]{cpp23standard} and local type inference~\cite{pierce2000local}: the type is synthesised from the body, and nothing flows into it from a use.
\Cpp{} refuses a use of the function before a return statement has deduced its type, and an FJlt body is a single return, so no body is typed through its own deduced type.
The one exception's self-call is such a use, and its function writes its result type.
A trait and its instances write \cppinline{auto} together or not at all, as the standard requires of a specialisation~\cite[{\texttt{dcl.spec.auto}}]{cpp23standard}.
A call's type arguments are deduced as well: \textsc{T-App}'s $\overline{\sigma}$, and a call $f\texttt{(}\overline{e}\texttt{)}$ is $f\texttt{<}\overline{\sigma}\texttt{>(}\overline{e}\texttt{)}$ at the $\overline{\sigma}$ the compiler deduces, so the rules check them and do not compute them.
A constant is typed under the empty contexts, and \textsc{T-Const} types its name at that type in every later declaration and in the program's term, unless a parameter or field of the same name is in $\Gamma$, which shadows it, as in \Cpp{}.
Its initialiser is a constant expression whose value keeps no allocation alive past the evaluation, as \cppinline{constexpr} requires of a variable and the compiler holds\cite[{\texttt{expr.const}}]{cpp23standard}.

\paragraph{Ranges and terms.}
A range $R\texttt{<}\tau\texttt{>}$ is an instance of a family $\mathcal{R}$ defined by conditions.
It has a length $|r| \in \mathbb{N}$, fixed before a fold starts, and, when $|r| > 0$, a head and a tail $\mathit{next}\,r$ with $|\mathit{next}\,r| = |r| - 1$.
\textsc{T-Range} types a range of $\mathcal{R}$ by its instance's conditions, \cppinline{iota(i, j)} only at $i \le j$, and \textsc{S-Fold} preserves them, since $\mathit{next}$ does.
\textsc{T-Range-Op} types \cppinline{size(r)} at \cppinline{size_t}, \cppinline{head(r)} at $\tau$ and \cppinline{next(r)} at $R\texttt{<}\tau\texttt{>}$.
The conditions are semantic, that a step descends or that two bounds are ordered.
A chain under a step that does not descend is therefore ill-typed, and the termination of a fold is, to that extent, assumed in the typing of its range.
A range may be a type variable $R$ under $\mathit{IsRange}\texttt{<}R\texttt{, }T\texttt{>}$ $=$ \cppinline{requires (const R r) { { size(r) } -> same_as<size_t>; { head(r) } -> same_as<T>; { next(r) } -> same_as<R>; }}.
The instances of $\mathcal{R}$ over $\tau$, and only they, satisfy it at $T = \tau$.
Membership in $\mathcal{R}$, with its falling length, is a premise, so three operations without that law do not suffice.
Listing~\ref{lst:rule3}'s \cppinline{concat} takes a range with a bounds-check, a count known before the fold starts, which the kernel reads as $\mathit{IsRange}$ with $T$ the deduced \cppinline{ranges::range_value_t<R>}.
A standard range with no such count, such as \cppinline{views::iota(a)}, is not in the kernel.
\textsc{T-Fold-R} types a fold at a type variable $R$ through the assumed shape $\mathit{IsRange}\texttt{<}R\texttt{, }\tau\texttt{>}$.
The length, head and tail are its three requirements, read from the record, and a caller's concrete range supplies the record by \textsc{E-Shape} at instantiation.
In \textsc{T-Cata}, $\mu F$ is a term type, the inductive type of the nodes that build it.
A term is a finite composition of nodes, and a node is a function object whose fields are terms.
The algebra is applied to each node together with the folded values of its legs, the node itself still in hand.
That is the library's interface (\cppinline{alg(t, cata(alg, leg)...)}), and the paramorphism's.

\paragraph{Well-founded declarations.}
Well-formedness also requires that the declarations, taken at their type arguments, be well founded under \emph{depends on}.
A shape at some types depends on the traits its requirements call and on its nested shapes, at those types.
A trait at some types depends on the context of every candidate whose head matches there, which resolution asks even of a candidate it then rejects, and on the candidate it picks.
An instance, a generic function or an operation at some types depends on the shapes of its context, on the declarations its body calls and on the operations its body constructs.
A constant depends on the declarations its initialiser names.
An operation also depends on the operations in its field, parameter and result types, each at the types it names.
The relation is on instantiations, not on bare declarations, since a structured instance may ask for the shape that asks for it, at a smaller type.
\cppinline{IsMonoid} at a pair calls \cppinline{unit} there, whose instance asks \cppinline{IsMonoid} at the components (Section~\ref{sec:translation:contexts}).
The measure that falls is the pair of the depth of the type arguments and the declaration's position.
Resolution, entailment, \textsc{E-Not}'s negative premise and \textsc{E-Law}'s evaluation all recurse along that relation.
Its being well founded makes the rules define one relation, and makes resolution terminate.
\Cpp{} has a rule of its own for a satisfaction that depends on itself, which the condition makes unnecessary.
Two syntactic conditions suffice, and the repository's linter reads the second.
A context on a structured head mentions only proper parts of its type arguments, so each step descends by the head; compare the Paterson conditions~\cite{sulzmann2007understanding}.
An instance whose head has no proper part at a slot, a type variable or a concrete type there, mentions in its context only shapes declared before the trait's primary.
A shape calling the trait is then never asked for while the trait is resolved.
A trait's specialisations may follow the shapes that call it, as Listing~\ref{lst:rule2}'s do, since a specialisation depends on the shape only through its body or its context.
A term type is required strictly positive: in the translated field types of its nodes, the type occurs only directly, never under the domain of a function type.
With one exception, a body calls only declarations it depends on.

\paragraph{The one exception.}
A function may call itself at the same instantiation on a value that visibly descends, Rule~3's step down.
The argument in the slot of a parameter $p$ is $p / c$ with $c \geq 2$, $p \gg c$ with $c \geq 1$, or $p - 1$.
It stands under a test of $p$ against $0$ whose branch returns without the call.
The parameter is of an unsigned carrier, or of a type variable under a shape the kernel reads as \cppinline{IsCountWord}.
That shape's fields are the operations the body uses and a coercion \cppinline{{ size_t(n) } -> same_as<size_t>} that preserves the count, so the count is no wider than \cppinline{size_t}.
This admits the binary method of Listing~\ref{lst:laws}, whose \cppinline{power} halves its count \cppinline{n} of type \cppinline{N}.
The library's \cppinline{cata} over typed terms is not this exception, as its recursion is on the type of the term, not on a value.
It is the kernel's primitive \cppinline{cata} spelled as a generic function, which is sugar (Appendix~\ref{app:fjlt:closure}).

\paragraph{Inferred laws.}
\label{app:fjlt:inferred}
A law at a construction is entailed, not declared.
The law is a trait call, resolution picks the instance whose context the program entails, and \textsc{E-Law} evaluates its certificate, which must be \texttt{true}.
Composed, these give a derived rule, the instance rule of the entailment of qualified types~\cite{jones1994qualified}:
\[
\textsc{(E-Law-Ctx)}\quad \dfrac{\begin{array}{c} \mathit{law}\texttt{<}\overline{\tau}\texttt{>} \text{ resolves to } \forall\overline{\alpha}.\,\overline{Q} \Rightarrow \mathit{law}\texttt{<}\overline{\tau'}\texttt{>} \qquad \overline{\tau} = \overline{\tau'}[\overline{\sigma}/\overline{\alpha}] \\ \Delta \Vdash \overline{Q}[\overline{\sigma}/\overline{\alpha}] \qquad \mathit{law}\texttt{<}\overline{\tau}\texttt{>} \Downarrow \texttt{true} \end{array}}{\Delta \Vdash \mathit{law}\texttt{<}\overline{\tau}\texttt{>}}
\]
A finite power of a semiring is a semimodule over it in this way.
The library declares the power a product of its base, and one instance, whose context is the base's semiring laws, gives the power's semimodule law, as semimodules over a semiring form a variety, closed under products.
Each step reads the laws of a term's immediate parts only, so the inference is local~\cite{pierce2000local}.
It ends by the conditions on contexts that end instance search~\cite{sulzmann2007understanding}.
